\documentclass{article}

% NeurIPS 2025 style
\usepackage[final]{neurips_2025}

\usepackage[utf8]{inputenc}
\usepackage[T1]{fontenc}
\usepackage{hyperref}
\usepackage{url}
\usepackage{booktabs}
\usepackage{amsfonts}
\usepackage{nicefrac}
\usepackage{microtype}
\usepackage{xcolor}
\usepackage{amsmath}
\usepackage{amssymb}
\usepackage{algorithm}
\usepackage{algorithmic}
\usepackage{graphicx}
\usepackage{subcaption}
\usepackage{multirow}
\usepackage{bm}

% Math operators
\DeclareMathOperator{\acosh}{acosh}
\DeclareMathOperator{\asinh}{asinh}
\DeclareMathOperator{\sign}{sign}
\DeclareMathOperator{\KL}{KL}
\DeclareMathOperator{\softmax}{softmax}
\DeclareMathOperator{\diag}{diag}
\newcommand{\R}{\mathbb{R}}
\newcommand{\Lo}{\mathbb{L}}
\newcommand{\E}{\mathbb{E}}
\newcommand{\mL}{\mathcal{L}}
\newcommand{\mH}{\mathcal{H}}
\newcommand{\mN}{\mathcal{N}}
\newcommand{\bx}{\bm{x}}
\newcommand{\bz}{\bm{z}}
\newcommand{\bmu}{\bm{\mu}}
\newcommand{\bsigma}{\bm{\sigma}}
\newcommand{\inner}[2]{\langle #1, #2 \rangle}

\title{Fine-Tuning Strategies for EEG Foundation Models:\\
A Comprehensive Technical Reference}

\author{
  EEG Foundation Model Project \\
  \texttt{NIPS\_finetune/}
}

\begin{document}

\maketitle

\begin{abstract}
This document provides a comprehensive technical reference for all fine-tuning methods implemented in the \texttt{NIPS\_finetune/} framework.
We present five distinct approaches to adapting pretrained EEG foundation models (CodeBrain SSSM and CBraMod) to downstream classification tasks:
(1)~\textbf{Standard Fine-Tuning} as baseline,
(2)~\textbf{Multi-Scale Fine-Tuning (MSFT)} with inter-layer cross-scale aggregation,
(3)~\textbf{Hyperbolic Fine-Tuning} with Lorentz manifold embeddings and domain-specific batch normalization,
(4)~\textbf{Mutual Information Fine-Tuning} combining variational information bottleneck with expert-guided InfoNCE,
and (5)~\textbf{Information Bottleneck + Disentanglement} with token-level compression and adversarial subject removal.
For each method, we provide precise mathematical formulations, algorithmic pseudocode, and architectural specifications.
Additionally, we describe a shared \textbf{Backbone Factory} with optional inter-layer bottleneck adapters for CBraMod.
\end{abstract}

%=============================================================================
\section{Preliminaries}
\label{sec:prelim}
%=============================================================================

\subsection{EEG Foundation Model Backbones}

All methods operate on a pretrained backbone $f_\theta$ that maps raw EEG segments to token-level representations.
The input is a 4D tensor $\bx \in \R^{B \times C \times S \times P}$ where $B$ is the batch size, $C$ is the number of EEG channels, $S$ is the number of temporal patches, and $P = 200$ is the patch size (samples per patch).
The backbone output is $\bm{H} = f_\theta(\bx) \in \R^{B \times C \times S \times D}$ with $D = 200$ for both backbones.

\paragraph{CodeBrain (SSSM).}
An S4-based state space model with $L = 8$ bidirectional residual blocks, diffusion-step embeddings, and optional VQ codebooks ($\text{codebook\_size} = 4096$).
The S4 sequence length is computed as $\text{s4\_lmax} = \lceil C \cdot S / 19 \rceil \times 19$.

\paragraph{CBraMod.}
A Criss-Cross Transformer with $L = 12$ layers, $d_\text{model} = 200$, $d_\text{ff} = 800$, $h = 8$ attention heads.
Each layer applies spatial attention (across channels) and temporal attention (across patches) in a criss-cross pattern.
The patch embedding module combines convolutional projection with spectral (FFT-based) features and learned depthwise-convolutional positional encoding.

\subsection{Notation}

\begin{table}[h]
\centering
\small
\begin{tabular}{@{}ll@{}}
\toprule
Symbol & Description \\
\midrule
$B, C, S, P, D$ & batch size, channels, temporal patches, patch size, feature dim \\
$T = C \times S$ & total number of tokens \\
$f_\theta$ & pretrained backbone (frozen unless noted) \\
$K < 0$ & negative curvature of hyperbolic space (default $-1$) \\
$\inner{\cdot}{\cdot}_\Lo$ & Minkowski inner product \\
$\oplus_K, \ominus_K, \otimes_K$ & Lorentz gyro-addition, inverse, scalar multiplication \\
$\mL_\text{CE}, \mL_\text{BCE}$ & cross-entropy, binary cross-entropy with logits \\
\bottomrule
\end{tabular}
\end{table}

%=============================================================================
\section{Method 1: Standard Fine-Tuning (Baseline)}
\label{sec:baseline}
%=============================================================================

\subsection{Architecture}

The baseline approach applies a multi-layer perceptron (MLP) classifier directly on the flattened backbone output.
The backbone may be either fully fine-tuned ($\texttt{--frozen False}$) or frozen ($\texttt{--frozen True}$, equivalent to linear probing).

\paragraph{Forward pass.}
\begin{equation}
\bm{H} = f_\theta(\bx) \in \R^{B \times C \times S \times D}, \quad
\hat{\bm{y}} = g_\phi\!\left(\text{flatten}(\bm{H})\right) \in \R^{B \times N_\text{cls}}
\end{equation}
where $g_\phi$ is the classifier head.

\paragraph{Classifier head} (default \texttt{all\_patch\_reps}, 3-layer MLP):
\begin{equation}
g_\phi(\bm{h}) = W_3\, \sigma\!\left( W_2\, \sigma\!\left( W_1 \bm{h} \right)\right)
\end{equation}
with $W_1 \in \R^{(S \cdot D) \times (C \cdot S \cdot D)}$, $W_2 \in \R^{D \times (S \cdot D)}$, $W_3 \in \R^{N_\text{cls} \times D}$, where $\sigma(\cdot) = \text{ELU}(\cdot)$ followed by dropout.

\subsection{Loss Function}

\begin{equation}
\mL = \mL_\text{CE}(\hat{\bm{y}}, \bm{y})
\end{equation}
For binary tasks, $\mL_\text{BCE}$ is used with sigmoid output (single logit).

\subsection{Pseudocode}

\begin{algorithm}[H]
\caption{Standard Fine-Tuning}
\begin{algorithmic}[1]
\REQUIRE Pretrained backbone $f_\theta$, classifier $g_\phi$, dataset $\mathcal{D}$
\FOR{each epoch}
  \FOR{each mini-batch $(\bx, \bm{y}) \sim \mathcal{D}$}
    \STATE $\bm{H} \leftarrow f_\theta(\bx)$ \COMMENT{Full backbone, optionally frozen}
    \STATE $\hat{\bm{y}} \leftarrow g_\phi(\text{flatten}(\bm{H}))$
    \STATE $\mL \leftarrow \mL_\text{CE}(\hat{\bm{y}}, \bm{y})$
    \STATE Update $\theta, \phi$ (or only $\phi$ if frozen) via $\nabla \mL$
  \ENDFOR
\ENDFOR
\end{algorithmic}
\end{algorithm}

%=============================================================================
\section{Method 2: Multi-Scale Fine-Tuning (MSFT)}
\label{sec:msft}
%=============================================================================

\subsection{Overview}
MSFT~\cite{qiao2025msft} injects \emph{trainable} cross-scale aggregation modules between the \emph{frozen} layers of the pretrained backbone.
The key idea is to generate multi-scale inputs via temporal downsampling and exchange information between scales at each layer, enabling the frozen backbone to process multi-resolution features.

\subsection{Multi-Scale Input Generation}

Given input $\bx \in \R^{B \times C \times S \times P}$, we generate $K+1$ scales ($k = 0, 1, \ldots, K$, default $K=2$ giving 3 scales) by recursive average pooling along the temporal dimension:
\begin{equation}
\bx^{(0)} = \bx, \quad
\bx^{(k)} = \text{AvgPool1D}\!\left(\bx^{(k-1)},\, \text{kernel}=2,\, \text{stride}=2\right), \quad
S_k = \left\lceil S / 2^k \right\rceil
\end{equation}

\subsection{Scale Adapter}

Each scale passes through a \emph{trainable} residual adapter after the frozen patch embedding:
\begin{equation}
\text{ScaleAdapter}(\bm{e}) = \bm{e} + \text{GELU}(W_\text{adapt}\, \bm{e} + b_\text{adapt}), \quad W_\text{adapt} \in \R^{D \times D}
\end{equation}

\subsection{Cross-Scale Aggregation}

At each frozen transformer/residual layer $\ell$, a trainable \texttt{CrossScaleAggregator} exchanges information between scales via bidirectional coarse-to-fine (C2F) and fine-to-coarse (F2C) pathways.

\paragraph{Coarse-to-Fine (C2F).}
For $k = K, K{-}1, \ldots, 1$:
\begin{equation}
\tilde{\bm{h}}^{(k-1)}_\text{c2f} = \bm{h}^{(k-1)} + \text{Upsample}\!\left(W^{(\ell)}_{\text{c2f},k}\, \bm{h}^{(k)},\; S_{k-1}\right)
\end{equation}
where $\text{Upsample}$ uses nearest-neighbor interpolation and $W^{(\ell)}_{\text{c2f},k} \in \R^{D \times D}$.

\paragraph{Fine-to-Coarse (F2C).}
For $k = 0, 1, \ldots, K{-}1$:
\begin{equation}
\tilde{\bm{h}}^{(k+1)}_\text{f2c} = \bm{h}^{(k+1)} + \text{AdaptiveAvgPool}\!\left(W^{(\ell)}_{\text{f2c},k}\, \bm{h}^{(k)},\; S_{k+1}\right)
\end{equation}

\paragraph{Merge.}
\begin{equation}
\bm{h}^{(k)}_\text{out} = \frac{1}{2}\!\left(\tilde{\bm{h}}^{(k)}_\text{c2f} + \tilde{\bm{h}}^{(k)}_\text{f2c}\right)
\end{equation}

\subsection{Per-Scale Classification and Mixing}

Each scale produces independent logits via a per-scale classifier $g_k$, combined with learnable softmax weights:
\begin{equation}
\hat{\bm{y}} = \sum_{k=0}^{K} w_k \cdot g_k(\bm{h}^{(k)}_\text{final}), \quad
\bm{w} = \softmax(\bm{\alpha}), \quad \bm{\alpha} \in \R^{K+1}
\end{equation}
where $\bm{\alpha}$ is a learnable parameter initialized to $\bm{0}$ (uniform mixing at init).

\subsection{Loss Function}

\begin{equation}
\mL = \mL_\text{CE}(\hat{\bm{y}}, \bm{y})
\end{equation}
No auxiliary losses; the contribution is purely architectural.

\subsection{Improved Variant: CrissCross-Aware MSFT}

The improved variant modifies the aggregation to preserve CBraMod's spatial-temporal attention structure:

\paragraph{Independent per-scale processing.} Each scale passes through each frozen transformer layer \emph{independently} (not concatenated), preserving the Criss-Cross Attention pattern.

\paragraph{Positional Refiner.} A depthwise Conv2D with BatchNorm is applied after each scale's patch embedding:
\begin{equation}
\text{PosRefine}(\bm{e}) = \bm{e} + \text{BN}\!\left(\text{DWConv2D}(\bm{e})\right), \quad \text{kernel} = 3 \times 3
\end{equation}

\paragraph{Spatial Fusion.} After temporal C2F/F2C, an additional per-scale spatial fusion:
\begin{equation}
\bm{h}^{(k)}_\text{fused} = \bm{h}^{(k)}_\text{agg} + \text{GELU}(W^{(k)}_\text{spatial}\, \bm{h}^{(k)}_\text{agg})
\end{equation}

\subsection{Pseudocode}

\begin{algorithm}[H]
\caption{Multi-Scale Fine-Tuning (MSFT)}
\begin{algorithmic}[1]
\REQUIRE Frozen backbone with layers $\{f^{(\ell)}\}_{\ell=1}^L$, $K{+}1$ scales
\STATE $\{\bx^{(k)}\}_{k=0}^K \leftarrow \text{MultiScaleGenerator}(\bx)$
\FOR{$k = 0$ to $K$}
  \STATE $\bm{e}^{(k)} \leftarrow \text{ScaleAdapter}_k\!\left(\text{FrozenPatchEmbed}(\bx^{(k)})\right)$
\ENDFOR
\STATE $\bm{h} \leftarrow \text{Concat}(\bm{e}^{(0)}, \ldots, \bm{e}^{(K)})$ along temporal dim
\FOR{$\ell = 1$ to $L$}
  \STATE $\bm{h} \leftarrow f^{(\ell)}(\bm{h})$ \COMMENT{Frozen layer}
  \STATE $\{\bm{h}^{(k)}\} \leftarrow \text{SplitByScale}(\bm{h})$
  \STATE $\{\bm{h}^{(k)}\} \leftarrow \text{CrossScaleAgg}^{(\ell)}(\{\bm{h}^{(k)}\})$ \COMMENT{Trainable}
  \STATE $\bm{h} \leftarrow \text{Concat}(\bm{h}^{(0)}, \ldots, \bm{h}^{(K)})$
\ENDFOR
\STATE $\{\bm{h}^{(k)}_\text{final}\} \leftarrow \text{SplitByScale}(\bm{h})$
\STATE $\bm{w} \leftarrow \softmax(\bm{\alpha})$
\STATE $\hat{\bm{y}} \leftarrow \sum_k w_k \cdot g_k(\bm{h}^{(k)}_\text{final})$
\STATE $\mL \leftarrow \mL_\text{CE}(\hat{\bm{y}}, \bm{y})$
\end{algorithmic}
\end{algorithm}

%=============================================================================
\section{Method 3: Hyperbolic Fine-Tuning}
\label{sec:hyperbolic}
%=============================================================================

\subsection{Overview}
Inspired by HEEGNet~\cite{li2026heegnet}, this method projects frozen backbone representations onto a Lorentz hyperbolic manifold, enabling hierarchical and domain-invariant feature learning.
The method combines \emph{hyperbolic multinomial logistic regression} (HMLR) with \emph{domain-specific hyperbolic batch normalization} (DSMDBN) and \emph{Horospherical Sliced-Wasserstein} (HHSW) regularization.

\subsection{Lorentz Manifold}

The Lorentz model of hyperbolic space is defined as:
\begin{equation}
\Lo^n_K = \left\{ \bx \in \R^{n+1} \;\middle|\; \inner{\bx}{\bx}_\Lo = \frac{1}{K},\; x_0 > 0 \right\}, \quad K < 0
\end{equation}
with the Minkowski inner product:
\begin{equation}
\inner{\bm{u}}{\bm{v}}_\Lo = -u_0 v_0 + \sum_{i=1}^{n} u_i v_i
\end{equation}

The origin is $\bm{o} = \left[\frac{1}{\sqrt{-K}},\, 0, \ldots, 0\right]^\top$.
The geodesic distance is:
\begin{equation}
d_K(\bx, \bm{y}) = \frac{1}{\sqrt{-K}} \acosh\!\left(K \inner{\bx}{\bm{y}}_\Lo\right)
\end{equation}

\paragraph{Exponential map at origin.}
\begin{equation}
\exp_{\bm{o}}(\bm{u}) = \left[\frac{\cosh(\vartheta)}{\sqrt{-K}},\; \frac{\sinh(\vartheta)}{\vartheta} \bm{u}_s \right], \quad
\vartheta = \sqrt{-K} \|\bm{u}_s\|_2
\end{equation}
where $\bm{u}_s = [u_1, \ldots, u_n]$ denotes the spatial components.

\paragraph{Projection to manifold.} Given Euclidean vector $\bm{v} \in \R^n$:
\begin{equation}
\text{proj}_K(\bm{v}) = \left[\sqrt{\|\bm{v}\|_2^2 - \frac{1}{K}},\; \bm{v}\right] \in \Lo^n_K
\end{equation}

\subsection{Architecture}

\paragraph{Hyperbolic Projection.}
\begin{equation}
\bm{h}_\text{flat} = \text{flatten}(f_\theta(\bx)) \in \R^{C \cdot S \cdot D}, \quad
\bm{v} = W_2\, \sigma(W_1\, \bm{h}_\text{flat}) \in \R^{n_\text{hyp}}
\end{equation}
\begin{equation}
\bm{p} = \text{proj}_K\!\left(\frac{\bm{v}}{\|\bm{v}\|_2}\right) \in \Lo^{n_\text{hyp}}_K
\end{equation}

\paragraph{Lorentz ELU Activation.}
\begin{equation}
\text{LorentzELU}(\bx) = \left[\sqrt{\|\text{ELU}(\bx_s)\|_2^2 - \frac{1}{K}},\; \text{ELU}(\bx_s)\right]
\end{equation}

\paragraph{Domain-Specific Momentum-then-Distribution BN (DSMDBN).}
For each domain $d$, maintain running Fr\'{e}chet mean $\bar{\bmu}_d$ and variance $\bar{\nu}^2_d$:
\begin{align}
\bmu_d^{(t)} &= \text{Fr\'{e}chetMean}\!\left(\{\bx_i : d_i = d\}\right) \\
\bar{\bmu}_d &\leftarrow \text{geodesic}\!\left(\bar{\bmu}_d,\, \bmu_d^{(t)},\, \eta\right) \\
\bar{\nu}^2_d &\leftarrow (1 - \eta)\, \bar{\nu}^2_d + \eta\, \nu^{2(t)}_d
\end{align}

Normalization:
\begin{equation}
\tilde{\bx} = \frac{s}{\sqrt{\bar{\nu}^2_d + \epsilon}} \otimes_K \left((\ominus_K \bar{\bmu}_d) \oplus_K \bx\right)
\end{equation}
where $s \in \R^+$ is a learnable scale parameter shared across domains.
The momentum $\eta$ decays during training: $\eta^{(k)} = \max(\eta_\text{min},\, \eta_0 \cdot \rho^k)$, with defaults $\eta_0 = 1.0$, $\rho = 0.99$, $\eta_\text{min} = 0.1$.

\paragraph{Hyperbolic Multinomial Logistic Regression (HMLR).}
For class $c$ with parameters $a_c \in \R$ (offset) and $\bz_c \in \R^{n}$ (normal vector):
\begin{align}
w_t^{(c)} &= \sinh(\sqrt{-K}\, a_c) \cdot \|\bz_c\|_2, \quad
\bm{w}_s^{(c)} = \cosh(\sqrt{-K}\, a_c) \cdot \bz_c \\
\beta_c &= \sqrt{-\left(w_t^{(c)}\right)^2 + \|\bm{w}_s^{(c)}\|_2^2} \\
\alpha_c(\bx) &= -w_t^{(c)}\, x_0 + \cosh(\sqrt{-K}\, a_c) \inner{\bx_s}{\bz_c} \\
d_c(\bx) &= \sqrt{|K|} \left| \asinh\!\left(\frac{\sqrt{-K}\, \alpha_c(\bx)}{\beta_c}\right) \right| \\
\text{logit}_c(\bx) &= \sign(\alpha_c(\bx)) \cdot \beta_c \cdot d_c(\bx)
\end{align}

\subsection{Loss Function}

\begin{equation}
\mL = \mL_\text{CE}(\hat{\bm{y}}, \bm{y}) + \lambda_\text{HHSW} \cdot \mL_\text{HHSW}
\end{equation}

\paragraph{HHSW (Horospherical Sliced-Wasserstein).}
For domain $d$ with features $\mathcal{P}_d = \{\bx_i : d_i = d\}$ and reference distribution $\mathcal{Q}$ (Gaussian mapped to $\Lo^n_K$ via $\exp_{\bm{o}}$):
\begin{equation}
\mL_\text{HHSW} = \sum_d \frac{1}{J} \sum_{j=1}^{J} W_2^2\!\left(\text{sort}(B_{\bm{v}_j}(\mathcal{P}_d)),\; \text{sort}(B_{\bm{v}_j}(\mathcal{Q}))\right)
\end{equation}
where $J = 1000$ random projections, $B_{\bm{v}}(\bz) = \log\!\left(-\inner{\bm{v} + \bm{o}}{\bz}_\Lo\right)$ is the Busemann function, and $W_2^2$ is the squared 1D Wasserstein-2 distance computed via sorted CDFs.

\subsection{Pseudocode}

\begin{algorithm}[H]
\caption{Hyperbolic Fine-Tuning}
\begin{algorithmic}[1]
\REQUIRE Frozen backbone $f_\theta$, modules $\{\text{HypProj}, \text{DSMDBN}, \text{HMLR}\}$
\FOR{each mini-batch $(\bx, \bm{y}, \bm{d})$}
  \STATE $\bm{H} \leftarrow f_\theta(\bx)$ \COMMENT{Frozen, no\_grad (unless adapters)}
  \STATE $\bm{h} \leftarrow \text{flatten}(\bm{H})$
  \STATE $\bm{p} \leftarrow \text{HypProj}(\bm{h})$ \COMMENT{$\R^d \to \Lo^n_K$}
  \STATE $\bm{p} \leftarrow \text{LorentzELU}(\bm{p})$
  \STATE $\tilde{\bm{p}} \leftarrow \text{DSMDBN}(\bm{p}, \bm{d})$ \COMMENT{Domain-specific BN}
  \STATE $\hat{\bm{y}} \leftarrow \text{HMLR}(\tilde{\bm{p}})$ \COMMENT{Hyperbolic classifier}
  \STATE $\mL \leftarrow \mL_\text{CE}(\hat{\bm{y}}, \bm{y}) + \lambda \cdot \mL_\text{HHSW}(\tilde{\bm{p}}, \bm{d})$
  \STATE Update trainable params via $\nabla \mL$ (RiemannianAdam)
  \STATE $\eta \leftarrow \max(\eta_\text{min},\, \eta \cdot \rho)$ \COMMENT{Momentum decay}
\ENDFOR
\end{algorithmic}
\end{algorithm}

\subsection{Optimizer}

Three parameter groups:
\begin{enumerate}
\item \textbf{Hyperbolic params} (HMLR, DSMDBN, LorentzELU): weight decay $= 0$
\item \textbf{Euclidean params} (HyperbolicProjection): weight decay $= 10^{-3}$
\item \textbf{Adapter params} (if CBraMod adapters): $\text{lr} = 0.1 \times \text{lr}_\text{base}$
\end{enumerate}
Optimizer: RiemannianAdam (via \texttt{geoopt}), fallback to Adam. Default $\text{lr} = 10^{-3}$.

%=============================================================================
\section{Method 4: MI Fine-Tuning (VIB + InfoNCE)}
\label{sec:mi}
%=============================================================================

\subsection{Overview}

This method combines two information-theoretic objectives with standard classification:
\begin{itemize}
\item \textbf{Variational Information Bottleneck (VIB):} Compress backbone representations to suppress subject-specific noise by minimizing $I(\bm{Z}; \bm{X})$ while maximizing $I(\bm{Z}; Y)$.
\item \textbf{InfoNCE:} Align backbone representations with expert-crafted EEG features (PSD, statistics) to maximize $I(\bm{Z}_\text{FM}; \bm{Z}_\text{expert})$.
\end{itemize}

\subsection{Architecture}

\paragraph{Representation Projection.}
\begin{equation}
\bm{Z}_\text{FM} = W_2\, \sigma(W_1\, \text{flatten}(\bm{H})) \in \R^{B \times H}, \quad H = 256
\end{equation}
where $\sigma = \text{GELU} \circ \text{Dropout}$.

\paragraph{VIB Layer.} Parameterizes $q(\bz | \bm{Z}_\text{FM}) = \mN(\bmu, \diag(\bsigma^2))$:
\begin{align}
\bmu &= f_\mu(\bm{Z}_\text{FM}) \in \R^{B \times V}, \quad
\log \bsigma^2 = \text{clamp}(f_{\log\sigma^2}(\bm{Z}_\text{FM}),\, {-10},\, 10) \\
\bz &= \bmu + \bsigma \odot \bm{\epsilon}, \quad \bm{\epsilon} \sim \mN(0, \bm{I}) \quad \text{(training)}
\end{align}
where $V = 128$ is the VIB dimension. In eval mode, $\bz = \bmu$.

\paragraph{Classifier Head.}
\begin{equation}
\hat{\bm{y}} = W_4\, \sigma(W_3\, \bz) \in \R^{B \times N_\text{cls}}
\end{equation}
with BatchNorm between layers.

\paragraph{Contrastive Head.}
\begin{equation}
\bz_\text{fm} = \frac{h_\text{contrast}(\bm{Z}_\text{FM})}{\|h_\text{contrast}(\bm{Z}_\text{FM})\|_2} \in \R^{B \times H}
\end{equation}

\paragraph{Expert Feature Extractor.}
Computes expert features $\bm{x}_\text{exp} \in \R^{B \times E}$ from raw EEG:
\begin{itemize}
\item \textbf{PSD}: FFT-based power in 5 bands ($\delta$: 0.5--4Hz, $\theta$: 4--8Hz, $\alpha$: 8--13Hz, $\beta$: 13--30Hz, $\gamma$: 30--50Hz), $\log(1 + \text{PSD})$ transformed. $E = 5C$.
\item \textbf{Statistics}: Mean, std, min, max per channel. $E = 4C$.
\item \textbf{Both}: Concatenation. $E = 9C$.
\end{itemize}

\paragraph{Expert Projector.}
\begin{equation}
\bz_\text{exp} = \frac{h_\text{expert}(\bm{x}_\text{exp})}{\|h_\text{expert}(\bm{x}_\text{exp})\|_2} \in \R^{B \times H}
\end{equation}

\subsection{Loss Function}

\begin{equation}
\mL = \mL_\text{CE}(\hat{\bm{y}}, \bm{y}) + \beta \cdot \mL_\text{VIB} + \alpha \cdot \mL_\text{InfoNCE}
\end{equation}

\paragraph{VIB Loss (KL divergence).}
\begin{equation}
\mL_\text{VIB} = \KL\!\left[q(\bz|\bx) \,\|\, p(\bz)\right] = -\frac{1}{2B} \sum_{i=1}^{B} \sum_{j=1}^{V} \left(1 + \log\sigma^2_{ij} - \mu^2_{ij} - \sigma^2_{ij}\right)
\end{equation}

\paragraph{Symmetric InfoNCE.}
\begin{equation}
\mL_\text{InfoNCE} = \frac{1}{2}\!\left[\mL_\text{NCE}(\bz_\text{fm} \to \bz_\text{exp}) + \mL_\text{NCE}(\bz_\text{exp} \to \bz_\text{fm})\right]
\end{equation}
where
\begin{equation}
\mL_\text{NCE}(\bm{A} \to \bm{B}) = -\frac{1}{B} \sum_{i=1}^{B} \log \frac{\exp(\bm{a}_i^\top \bm{b}_i / \tau)}{\sum_{j=1}^{B} \exp(\bm{a}_i^\top \bm{b}_j / \tau)}
\end{equation}
with temperature $\tau = 0.07$.

\subsection{Pseudocode}

\begin{algorithm}[H]
\caption{MI Fine-Tuning (VIB + InfoNCE)}
\begin{algorithmic}[1]
\REQUIRE Frozen backbone $f_\theta$, MIFineTuner, ExpertExtractor
\FOR{each mini-batch $(\bx, \bm{y})$}
  \STATE $\bm{H} \leftarrow f_\theta(\bx)$ \COMMENT{Frozen backbone}
  \STATE $\bm{Z}_\text{FM} \leftarrow \text{RepProj}(\text{flatten}(\bm{H}))$
  \STATE $\bz, \bmu, \log\bsigma^2 \leftarrow \text{VIB}(\bm{Z}_\text{FM})$ \COMMENT{Reparameterize}
  \STATE $\hat{\bm{y}} \leftarrow \text{Classifier}(\bz)$
  \STATE $\bz_\text{fm} \leftarrow \ell_2\text{-norm}(\text{ContrastHead}(\bm{Z}_\text{FM}))$
  \STATE $\bm{x}_\text{exp} \leftarrow \text{ExpertExtract}(\bx)$ \COMMENT{PSD / stats}
  \STATE $\bz_\text{exp} \leftarrow \ell_2\text{-norm}(\text{ExpertProj}(\bm{x}_\text{exp}))$
  \STATE $\mL \leftarrow \mL_\text{CE} + \beta \cdot \mL_\text{VIB}(\bmu, \log\bsigma^2) + \alpha \cdot \mL_\text{InfoNCE}(\bz_\text{fm}, \bz_\text{exp})$
  \STATE Update trainable params via $\nabla \mL$
\ENDFOR
\end{algorithmic}
\end{algorithm}

\subsection{Default Hyperparameters}

\begin{table}[h]
\centering
\small
\begin{tabular}{@{}lll@{}}
\toprule
Parameter & Symbol & Default \\
\midrule
Hidden dimension & $H$ & 256 \\
VIB dimension & $V$ & 128 \\
InfoNCE weight & $\alpha$ & 1.0 \\
VIB weight & $\beta$ & $10^{-3}$ \\
Temperature & $\tau$ & 0.07 \\
Expert features & -- & PSD (5 bands) \\
Optimizer & -- & AdamW \\
Learning rate & -- & $10^{-3}$ \\
\bottomrule
\end{tabular}
\end{table}

%=============================================================================
\section{Method 5: Information Bottleneck + Disentanglement}
\label{sec:ib}
%=============================================================================

\subsection{Overview}

This method extends the information-theoretic approach with two key innovations over Method~4:
\begin{enumerate}
\item \textbf{Token-level IB}: Applies the information bottleneck at the per-token level (preserving spatial-temporal structure) rather than after global flattening.
\item \textbf{Gradient Reversal Layer (GRL)}: Explicitly minimizes $I(\bm{Z}; \text{Subject})$ via adversarial training, eliminating the need for hand-crafted expert features.
\end{enumerate}

\subsection{Architecture}

\paragraph{Token-Level IB Adapter.}
The backbone output $\bm{H} \in \R^{B \times C \times S \times D}$ is reshaped to $\bm{H} \in \R^{B \times T \times D}$ with $T = C \times S$.
A shared encoder produces per-token variational parameters:
\begin{align}
\bm{e}_t &= \sigma(\text{LN}(W_\text{enc}\, \bm{h}_t)) \in \R^{D}, \quad \forall t \in \{1, \ldots, T\} \\
\bmu_t &= W_\mu\, \bm{e}_t \in \R^{D'}, \quad
\log\bsigma^2_t = \text{clamp}(W_{\log\sigma^2}\, \bm{e}_t,\, {-10},\, 10) \\
\bz_t &= \bmu_t + \bsigma_t \odot \bm{\epsilon}_t, \quad \bm{\epsilon}_t \sim \mN(0, \bm{I})
\end{align}
where $D' = 128$ (latent dimension), $\sigma = \text{GELU} \circ \text{Dropout}$, and LN denotes LayerNorm.

\paragraph{Aggregation.}
\begin{equation}
\bz_\text{agg} = \frac{1}{T} \sum_{t=1}^{T} \bz_t \in \R^{D'}
\end{equation}

\paragraph{Disease Classifier.}
\begin{equation}
\hat{\bm{y}}_\text{disease} = g_\text{disease}(\bz_\text{agg}) \in \R^{N_\text{cls}}
\end{equation}

\paragraph{Gradient Reversal Layer.}
\begin{equation}
\text{GRL}_\lambda(\bz) \triangleq \bz, \quad
\frac{\partial \text{GRL}_\lambda}{\partial \bz} = -\lambda \cdot \bm{I}
\end{equation}
Identity in forward pass, negated and scaled gradients in backward pass.

\paragraph{Subject Classifier.}
\begin{equation}
\hat{\bm{y}}_\text{subject} = g_\text{subject}(\text{GRL}_\lambda(\bz_\text{agg})) \in \R^{N_\text{subj}}
\end{equation}

\subsection{Loss Function}

\begin{equation}
\mL = \mL_\text{task} + \beta \cdot \mL_\text{IB} + \lambda_\text{adv} \cdot \mL_\text{MI}
\end{equation}

\paragraph{Task Loss.}
$\mL_\text{task} = \mL_\text{CE}(\hat{\bm{y}}_\text{disease}, \bm{y})$ or $\mL_\text{BCE}$ for binary tasks.

\paragraph{Token-Level IB Loss.}
\begin{equation}
\mL_\text{IB} = -\frac{1}{2BT} \sum_{i=1}^{B} \sum_{t=1}^{T} \sum_{j=1}^{D'} \left(1 + \log\sigma^2_{i,t,j} - \mu^2_{i,t,j} - \sigma^2_{i,t,j}\right)
\end{equation}
Compared to Method~4's VIB (averaged over $B$), this averages over $B \times T$, giving per-token compression.

\paragraph{Adversarial MI Loss.}
\begin{equation}
\mL_\text{MI} = \mL_\text{CE}(\hat{\bm{y}}_\text{subject}, \bm{s})
\end{equation}
where $\bm{s}$ are subject labels. The GRL ensures that minimizing $\mL_\text{MI}$ w.r.t.\ the subject head simultaneously \emph{maximizes} $\mL_\text{MI}$ w.r.t.\ the encoder and IB adapter --- no alternating min-max required.

\subsection{GRL Lambda Schedule}

The adversarial pressure ramps up during training via a sigmoid schedule:
\begin{equation}
\lambda(p) = \frac{2}{1 + \exp(-\gamma \cdot p)} - 1, \quad p = \frac{\text{epoch}}{\text{total\_epochs}}
\end{equation}
with $\gamma = 10$, ramping from $\lambda \approx 0$ to $\lambda \approx 1$.

\subsection{Interpretability: Information Retention Heatmap}

The per-token KL divergence provides a natural measure of \emph{how much information} each spatial-temporal token retains after compression:
\begin{equation}
\text{KL}_t = -\frac{1}{2} \sum_{j=1}^{D'} \left(1 + \log\sigma^2_{t,j} - \mu^2_{t,j} - \sigma^2_{t,j}\right)
\end{equation}
Reshaped to $(C, S)$, this yields a spatial-temporal heatmap.
Channel importance: $\text{CI}_c = \frac{1}{S}\sum_s \text{KL}_{c,s}$.
Temporal importance: $\text{TI}_s = \frac{1}{C}\sum_c \text{KL}_{c,s}$.

\subsection{Pseudocode}

\begin{algorithm}[H]
\caption{IB + Disentanglement Fine-Tuning}
\begin{algorithmic}[1]
\REQUIRE Frozen backbone $f_\theta$, IB Adapter, Disease Head, GRL, Subject Head
\FOR{each epoch $e$}
  \STATE $\lambda \leftarrow 2/(1 + \exp(-\gamma \cdot e / E)) - 1$ \COMMENT{Update GRL}
  \FOR{each mini-batch $(\bx, \bm{y}, \bm{s})$}
    \STATE $\bm{H} \leftarrow \text{reshape}(f_\theta(\bx),\, [B, T, D])$ \COMMENT{Frozen backbone}
    \STATE $\bm{Z}, \bmu, \log\bsigma^2 \leftarrow \text{IB\_Adapter}(\bm{H})$ \COMMENT{Token-level}
    \STATE $\bz_\text{agg} \leftarrow \text{mean}(\bm{Z}, \text{dim}=1)$
    \STATE $\hat{\bm{y}} \leftarrow g_\text{disease}(\bz_\text{agg})$
    \STATE $\hat{\bm{s}} \leftarrow g_\text{subject}(\text{GRL}_\lambda(\bz_\text{agg}))$
    \STATE $\mL \leftarrow \mL_\text{CE}(\hat{\bm{y}}, \bm{y}) + \beta \cdot \mL_\text{IB}(\bmu, \log\bsigma^2) + \lambda_\text{adv} \cdot \mL_\text{CE}(\hat{\bm{s}}, \bm{s})$
    \STATE Clip gradients: $\|\nabla\| \leq 5.0$
    \STATE Update trainable params via $\nabla \mL$
  \ENDFOR
\ENDFOR
\end{algorithmic}
\end{algorithm}

\subsection{Optimizer}

Four parameter groups with AdamW:
\begin{enumerate}
\item \textbf{IB Adapter}: $\text{lr} = \text{lr}_\text{base}$
\item \textbf{Disease Head}: $\text{lr} = \text{lr}_\text{base}$
\item \textbf{Subject Head}: $\text{lr} = 2 \times \text{lr}_\text{base}$ (faster adversarial learning)
\item \textbf{Backbone Adapters} (if present): $\text{lr} = 0.1 \times \text{lr}_\text{base}$
\end{enumerate}

%=============================================================================
\section{Shared Infrastructure: Backbone Factory with Adapters}
\label{sec:factory}
%=============================================================================

\subsection{Bottleneck Adapter}

For CBraMod with inter-layer adapter injection, we insert a lightweight residual bottleneck after each frozen transformer layer:
\begin{equation}
\text{Adapter}(\bm{h}) = \bm{h} + W_\text{up}\, \text{GELU}(W_\text{down}\, \bm{h})
\end{equation}
where $W_\text{down} \in \R^{(D/r) \times D}$, $W_\text{up} \in \R^{D \times (D/r)}$ with reduction factor $r = 4$ (default), giving $D/r = 50$.

\paragraph{Zero initialization.} $W_\text{up}$ and its bias are initialized to zero, so the adapter acts as an identity function at initialization, preserving pretrained features.

\paragraph{Parameters per adapter:} $D \cdot (D/r) + (D/r) + (D/r) \cdot D + D = 2D(D/r) + D + D/r \approx 20{,}250$ for $D=200, r=4$.
With 12 layers: $\sim$243K total adapter parameters.

\subsection{Forward Pass Comparison}

\paragraph{HEAD-ONLY (default).}
\begin{equation}
\bx \xrightarrow{\text{frozen}} f_\theta(\bx) \xrightarrow{\text{trainable}} \text{Head}(\cdot) \to \hat{\bm{y}}
\end{equation}

\paragraph{INTER-LAYER (with adapters).}
\begin{equation}
\bx \xrightarrow{\text{frozen PE}} \bm{e} \xrightarrow[\text{adapter}_1]{\text{frozen layer}_1} \cdots \xrightarrow[\text{adapter}_L]{\text{frozen layer}_L} \bm{h} \xrightarrow{\text{trainable}} \text{Head}(\cdot) \to \hat{\bm{y}}
\end{equation}

%=============================================================================
\section{Ablation Configurations}
\label{sec:ablation}
%=============================================================================

\subsection{IB + Disentanglement Ablations}

\begin{table}[h]
\centering
\small
\begin{tabular}{@{}llll@{}}
\toprule
Configuration & $\beta$ (IB) & $\lambda_\text{adv}$ (GRL) & Description \\
\midrule
\texttt{full\_ib\_adv} & $10^{-3}$ & 0.5 & Full model \\
\texttt{ib\_only} & $10^{-3}$ & 0.0 & IB without adversarial \\
\texttt{adv\_only} & 0.0 & 0.5 & Adversarial without IB \\
\texttt{baseline\_ce} & 0.0 & 0.0 & CE only (head-only baseline) \\
\texttt{high\_beta} & $10^{-2}$ & 0.5 & Stronger compression \\
\texttt{low\_beta} & $10^{-4}$ & 0.5 & Weaker compression \\
\texttt{high\_lambda} & $10^{-3}$ & 1.0 & Stronger disentanglement \\
\texttt{low\_lambda} & $10^{-3}$ & 0.1 & Weaker disentanglement \\
\bottomrule
\end{tabular}
\end{table}

\subsection{MSFT Scale Ablations}

Number of scales $K+1 \in \{1, 2, 3, 4\}$, tested across datasets (TUEV, TUAB).

%=============================================================================
\section{Comparison Summary}
\label{sec:comparison}
%=============================================================================

\begin{table}[h]
\centering
\small
\setlength{\tabcolsep}{3pt}
\begin{tabular}{@{}p{2.0cm}p{3.8cm}p{2.5cm}p{2.2cm}p{2.5cm}@{}}
\toprule
Method & Loss & Trainable Modules & Backbone Mode & Key Innovation \\
\midrule
Baseline & $\mL_\text{CE}$ & Backbone + Classifier & Full FT / Frozen & -- \\
\addlinespace
MSFT & $\mL_\text{CE}$ & ScaleAdapters, CrossScaleAggs, Classifiers & Frozen (inter-layer) & Multi-scale C2F/F2C aggregation \\
\addlinespace
Hyperbolic & $\mL_\text{CE} + \lambda \mL_\text{HHSW}$ & HypProj, DSMDBN, HMLR & Frozen (head-only) & Lorentz manifold, domain BN \\
\addlinespace
MI (VIB+NCE) & $\mL_\text{CE} + \beta \mL_\text{VIB} + \alpha \mL_\text{NCE}$ & RepProj, VIB, Contrast, ExpertProj, Classifier & Frozen (head-only) & Expert-guided contrastive + VIB \\
\addlinespace
IB+Disentangle & $\mL_\text{task} + \beta \mL_\text{IB} + \lambda \mL_\text{MI}$ & IB Adapter, Disease Head, Subject Head & Frozen (head-only) & Token-level IB, GRL adversarial \\
\bottomrule
\end{tabular}
\caption{Comparison of all fine-tuning methods. All methods except baseline keep the backbone frozen. Methods 3--5 support optional CBraMod inter-layer adapters via the Backbone Factory.}
\end{table}

%=============================================================================
\section{Supported Backbones and Datasets}
\label{sec:support}
%=============================================================================

\begin{table}[h]
\centering
\small
\begin{tabular}{@{}lccccc@{}}
\toprule
 & Baseline & MSFT & Hyperbolic & MI & IB+Disentangle \\
\midrule
CodeBrain (SSSM) & \checkmark & \checkmark & \checkmark & \checkmark & \checkmark \\
CBraMod & \checkmark & \checkmark & \checkmark & \checkmark & \checkmark \\
CBraMod + Adapters & -- & -- & \checkmark & \checkmark & \checkmark \\
\bottomrule
\end{tabular}
\caption{Backbone support matrix across methods.}
\end{table}

\begin{table}[h]
\centering
\small
\begin{tabular}{@{}lcccccc@{}}
\toprule
Dataset & Task & Classes & Channels & Sampling & Duration & Patch Size \\
\midrule
TUEV & Multiclass & 6 & 16 & 200\,Hz & 5\,s & 200 \\
TUAB & Binary & 2 & 16 & 200\,Hz & 10\,s & 200 \\
CHB-MIT & Binary & 2 & 16 & 256\,Hz & 5\,s & 200 \\
TUSZ & Binary & 2 & 22 & 200\,Hz & 5\,s & 200 \\
DIAGNOSIS & Multiclass & 4 & 58 & 250\,Hz & 1\,s & 200 \\
DEPRESSION & Binary & 2 & 63 & 200\,Hz & 5\,s & 200 \\
\bottomrule
\end{tabular}
\caption{Supported EEG datasets with configurations.}
\end{table}

\begin{thebibliography}{9}
\bibitem{qiao2025msft}
Qiao et al.
\newblock Multi-Scale Finetuning for Encoder-based Time Series Foundation Models.
\newblock 2025.

\bibitem{li2026heegnet}
Li et al.
\newblock HEEGNet: Hyperbolic EEG Foundation Model.
\newblock 2026.

\bibitem{alemi2017deep}
Alemi, A.~A., Fischer, I., Dillon, J.~V., and Murphy, K.
\newblock Deep variational information bottleneck.
\newblock In \emph{ICLR}, 2017.

\bibitem{ganin2016domain}
Ganin, Y. et al.
\newblock Domain-adversarial training of neural networks.
\newblock \emph{JMLR}, 17(1):2096--2030, 2016.

\bibitem{houlsby2019parameter}
Houlsby, N. et al.
\newblock Parameter-efficient transfer learning for NLP.
\newblock In \emph{ICML}, 2019.

\bibitem{oord2018representation}
van den Oord, A., Li, Y., and Vinyals, O.
\newblock Representation learning with contrastive predictive coding.
\newblock \emph{arXiv preprint arXiv:1807.03748}, 2018.
\end{thebibliography}

\end{document}
